% ================================================================
%  Situation Réelle 1 — Réhabilitation thermique d'une maison individuelle
%  BTS FED — Option GCF — Semestre 1
%  Chapitres mobilisés : 01 (thermique), 02 (hydraulique), 04 (production), 05 (émission)
% ================================================================

\chapter{SR-1 — Réhabilitation thermique d'une maison individuelle}

\begin{tcolorbox}[
  enhanced, colback=FedGreen!8, colframe=FedGreen, arc=3pt,
  fonttitle=\bfseries\sffamily\large, coltitle=white, colbacktitle=FedGreen,
  title={Situation Réelle 1 — Fiche d'identité du projet},
  width=\textwidth
]
\begin{tabular}{@{}ll@{\hspace{2cm}}ll@{}}
  \textbf{Semestre} & 1 &
  \textbf{Durée} & 6 semaines \\[4pt]
  \textbf{Groupe} & 2 à 3 étudiants &
  \textbf{Chapitres} & 01, 02, 04, 05 \\[4pt]
  \textbf{Note 1} & Dossier technique (coeff.\ 3) &
  \textbf{Note 2} & Soutenance orale (coeff.\ 2) \\[4pt]
  \textbf{Compétences} & C2-3, C3-3, C5-1, C8-3, C8-4 & & \\
\end{tabular}
\end{tcolorbox}

\bigskip

% ================================================================
\section{Contexte professionnel}
% ================================================================

Vous êtes technicien supérieur GCF dans l'entreprise \textbf{ThermoPro SARL} (Grenoble, 38).
Un particulier, M.\ et Mme Aubert, souhaite rénover leur maison individuelle construite
en \textbf{1975} afin d'améliorer son étiquette DPE de \textbf{D} vers \textbf{B} et
de remplacer la chaudière fioul existante par une \textbf{pompe à chaleur air/eau}.

\subsection*{Description du bâtiment}

\begin{figure}[htbp]
  \centering
  \begin{tikzpicture}[scale=1.0, >=Stealth]

    % --- Sol ---
    \fill[gray!20] (-0.3,-0.3) rectangle (10.3, 0);
    \draw[thick] (-0.3,0) -- (10.3,0);

    % --- Murs extérieurs ---
    \draw[thick, FedBlue, fill=FedBlue!10]
      (0,0) rectangle (10,4.5);

    % --- Toit ---
    \draw[thick, FedBlue!80, fill=FedBlue!20]
      (-0.5,4.5) -- (5,7) -- (10.5,4.5) -- cycle;

    % --- Plancher intermédiaire (R+1) ---
    \draw[FedOrange, thick] (0,2.2) -- (10,2.2);
    \node[FedOrange, font=\footnotesize] at (5,2.0) {Plancher intermédiaire};

    % --- Fenêtres RDC ---
    \draw[FedBlue!60, fill=cyan!15, thick] (1,0.5) rectangle (2.5,1.8);
    \draw[FedBlue!60, fill=cyan!15, thick] (4,0.5) rectangle (5.8,1.8);
    \draw[FedBlue!60, fill=cyan!15, thick] (7.2,0.5) rectangle (8.8,1.8);
    % --- Porte ---
    \draw[brown!70, fill=brown!20, thick] (3.1,0) rectangle (3.9,2.0);
    \draw[brown!70] (3.1,1.0) arc (180:0:0.4);

    % --- Fenêtres R+1 ---
    \draw[FedBlue!60, fill=cyan!15, thick] (1.5,2.6) rectangle (3,3.8);
    \draw[FedBlue!60, fill=cyan!15, thick] (5,2.6) rectangle (6.5,3.8);
    \draw[FedBlue!60, fill=cyan!15, thick] (7.5,2.6) rectangle (9,3.8);

    % --- Velux ---
    \draw[FedBlue!60, fill=cyan!20, thick] (3.5,5.2) rectangle (4.8,6.0);
    \draw[FedBlue!60, fill=cyan!20, thick] (5.5,5.2) rectangle (6.8,6.0);

    % --- Annotations ---
    \node[font=\small\bfseries, FedBlue] at (5, 1.0) {RDC — Surface : \SI{60}{\metre\squared}};
    \node[font=\small\bfseries, FedBlue] at (5, 3.3) {R+1 — Surface : \SI{60}{\metre\squared}};
    \node[font=\small, FedBlue!80] at (5, 5.7) {Combles perdus};

    % --- Côtes ---
    \draw[<->, gray] (-0.1, 0) -- (-0.1, 4.5)
      node[midway, left, font=\footnotesize] {\SI{5.5}{\metre}};
    \draw[<->, gray] (0, -0.2) -- (10, -0.2)
      node[midway, below, font=\footnotesize] {\SI{10}{\metre}};

    % --- Isolation existante (hachures) ---
    \foreach \y in {0.3, 0.7, 1.1, 1.5, 1.9, 2.4, 2.8, 3.2, 3.6, 4.0, 4.3}{
      \draw[red!40, thin] (0,\y) -- (0.25, \y+0.2);
      \draw[red!40, thin] (9.75,\y) -- (10,\y+0.2);
    }
    \node[red!60, font=\tiny, rotate=90] at (0.15, 2.5) {Laine verre 60mm};

    % --- Chaudière fioul existante ---
    \draw[FedRed, fill=FedRed!15, thick, rounded corners=2pt]
      (0.5, -2.5) rectangle (2.5, -1.0);
    \node[font=\footnotesize\bfseries, FedRed] at (1.5,-1.75) {Chaudière fioul};
    \node[font=\tiny, FedRed] at (1.5,-2.1) {18 kW — η = 82\%};
    \draw[FedRed, thick, ->] (1.5,-1.0) -- (1.5,-0.0);

    % --- PAC future ---
    \draw[FedGreen, fill=FedGreen!15, thick, rounded corners=2pt, dashed]
      (7, -2.5) rectangle (9.5, -1.0);
    \node[font=\footnotesize\bfseries, FedGreen] at (8.25,-1.6) {PAC air/eau};
    \node[font=\tiny, FedGreen] at (8.25,-2.1) {à dimensionner};
    \node[font=\tiny, FedGreen] at (8.25,-2.35) {A7/W35};

    % --- Légende ---
    \draw[FedRed, thick] (0.5,-3.2) -- (1.0,-3.2);
    \node[right, font=\tiny] at (1.0,-3.2) {Existant (dépose)};
    \draw[FedGreen, thick, dashed] (3.5,-3.2) -- (4.0,-3.2);
    \node[right, font=\tiny] at (4.0,-3.2) {Projet};
    \draw[red!40] (6.5,-3.2) -- (7.0,-3.2);
    \node[right, font=\tiny] at (7.0,-3.2) {Isolation existante};

  \end{tikzpicture}
  \caption{Coupe schématique de la maison Aubert — Grenoble (H1c) — Construction 1975
    — S\textsubscript{habitable} = \SI{120}{\metre\squared} — Chaudière fioul à remplacer par PAC air/eau}
  \label{fig:sr1-coupe-maison}
\end{figure}

\subsection*{Données du bâtiment existant}

\begin{center}
\begin{tabular}{lllr}
  \toprule
  \textbf{Élément} & \textbf{Composition actuelle} & \textbf{U actuel} & \textbf{Proportion} \\
  \midrule
  Murs extérieurs & Parpaings 20 cm + enduits & \SI{1.8}{\watt\per\metre\squared\per\kelvin} & \SI{140}{\metre\squared} \\
  Plancher bas & Dalle béton sur terre-plein & \SI{0.6}{\watt\per\metre\squared\per\kelvin} & \SI{60}{\metre\squared} \\
  Toiture & Tuiles + laine de verre 60 mm & \SI{0.55}{\watt\per\metre\squared\per\kelvin} & \SI{65}{\metre\squared} \\
  Fenêtres RDC & Simple vitrage bois & \SI{4.5}{\watt\per\metre\squared\per\kelvin} & \SI{9}{\metre\squared} \\
  Fenêtres R+1 & Simple vitrage bois & \SI{4.5}{\watt\per\metre\squared\per\kelvin} & \SI{9}{\metre\squared} \\
  Porte d'entrée & Bois plein & \SI{3.2}{\watt\per\metre\squared\per\kelvin} & \SI{2}{\metre\squared} \\
  \midrule
  \textbf{Pont thermique linéique} & Jonction mur/plancher & $\psi = \SI{0.8}{\watt\per\metre\per\kelvin}$ & \SI{40}{\metre} \\
  \bottomrule
\end{tabular}
\end{center}

\begin{definition}[Données climatiques — Grenoble, zone H1c]
  Température extérieure de base : $T_{ext,base} = \SI{-11}{\celsius}$ (NF EN 12831) \\
  Température intérieure de confort : $T_{int} = \SI{19}{\celsius}$ \\
  Degré-jours unifiés (DJU) : \textbf{2 900 DJU/an} (référence 18°C) \\
  Ensoleillement annuel : \SI{1950}{\hour\per\year}
\end{definition}

% ================================================================
\section{Schéma de principe de l'installation projetée}
% ================================================================

\begin{figure}[htbp]
  \centering
  \begin{tikzpicture}[
    >=Stealth, thick,
    bloc/.style={draw, rounded corners=3pt, minimum width=2.8cm, minimum height=1.0cm,
                 align=center, font=\small},
    tuyau/.style={draw, ->, thick},
    dep/.style={FedOrange, thick},   % départ chaud
    ret/.style={FedBlue, thick},     % retour froid
  ]

    % === PAC air/eau ===
    \node[bloc, fill=FedGreen!15, draw=FedGreen] (pac) at (0,0)
      {PAC air/eau\\{\footnotesize A7/W35}};

    % === Ballon tampon ===
    \node[bloc, fill=FedBlue!10, draw=FedBlue] (tampon) at (4,0)
      {Ballon\\tampon\\{\footnotesize 300 L}};

    % === Collecteur départ ===
    \node[bloc, fill=FedOrange!15, draw=FedOrange, minimum width=3.5cm] (coll) at (8,0)
      {Collecteur\\départ/retour};

    % === Émetteurs ===
    \node[bloc, fill=gray!10, draw=gray] (pch-rdc) at (11, 2.5)
      {PCH RDC\\{\footnotesize 60 m²}};
    \node[bloc, fill=gray!10, draw=gray] (rad-r1) at (11, -2.5)
      {Radiateurs R+1\\{\footnotesize BT 45/40°C}};

    % === Vase d'expansion ===
    \node[draw, circle, minimum size=0.8cm, fill=yellow!20, font=\tiny] (vase) at (4, 2.5)
      {VE};

    % === Pompes de circulation ===
    \node[draw, circle, minimum size=0.6cm, fill=gray!20, font=\tiny] (pompe1) at (2, 0.4)
      {P1};
    \node[draw, circle, minimum size=0.6cm, fill=gray!20, font=\tiny] (pompe2) at (6, 0.4)
      {P2};

    % === Connexions départ (orange) ===
    \draw[dep, ->] (pac.east) -- node[above, font=\tiny] {55°C} (tampon.west);
    \draw[dep, ->] (tampon.east) -- node[above, font=\tiny] {45°C} (coll.west);
    \draw[dep, ->] (coll.north east) |- (pch-rdc.west);
    \draw[dep, ->] (coll.south east) |- (rad-r1.west);

    % === Connexions retour (bleu) ===
    \draw[ret, ->] (pch-rdc.south) -- ++(0,-0.8) -| (coll.north);
    \draw[ret, ->] (rad-r1.north) -- ++(0,0.8) -| (coll.south);
    \draw[ret, ->] (coll.west|-tampon) -- node[below, font=\tiny] {35°C} (tampon.east|-tampon);
    \draw[ret, ->] (tampon.west) -- node[below, font=\tiny] {40°C} (pac.east|-pac.south);

    % === Vase d'expansion connexion ===
    \draw[gray, dashed] (vase) -- (tampon.north);

    % === Légende ===
    \draw[dep, ->] (0,-4.2) -- (1.2,-4.2) node[right, font=\tiny, black] {Circuit départ chaud};
    \draw[ret, ->] (3.5,-4.2) -- (4.7,-4.2) node[right, font=\tiny, black] {Circuit retour froid};
    \node[draw, circle, minimum size=0.5cm, fill=gray!20, font=\tiny] at (7.5,-4.2) {P};
    \node[right, font=\tiny] at (7.8,-4.2) {Circulateur};
    \node[draw, circle, minimum size=0.5cm, fill=yellow!20, font=\tiny] at (9.5,-4.2) {VE};
    \node[right, font=\tiny] at (9.8,-4.2) {Vase expansion};

    % === Sonde extérieure ===
    \node[draw, diamond, fill=purple!15, minimum size=0.5cm, font=\tiny] (sonde) at (-2, 2.0)
      {$T_e$};
    \draw[purple, dashed, ->] (sonde) -- (pac.north west);
    \node[purple, font=\tiny, above] at (-1, 2.0) {Régulation};
    \node[purple, font=\tiny, above] at (-1, 1.6) {climatique};

  \end{tikzpicture}
  \caption{Schéma hydraulique de principe — Installation PAC air/eau + ballon tampon + PCH RDC + radiateurs R+1
    — Régulation climatique par sonde extérieure}
  \label{fig:sr1-schema-hydraulique}
\end{figure}

% ================================================================
\section{Cahier des charges}
% ================================================================

\begin{methode}[Livrables attendus — Note 1 : Dossier technique écrit]
\textbf{À remettre en semaine 6 — Format PDF, 25 pages minimum}

\begin{enumerate}
  \item \textbf{Bilan thermique} selon NF EN 12831 :
    \begin{itemize}
      \item Calcul du coefficient $U$ de chaque paroi (formule de la résistance thermique équivalente)
      \item Calcul des déperditions totales $\Phi_{dep}$ du bâtiment [\si{\watt}]
      \item Calcul de $U_{bat}$ et comparaison avec le seuil RE2020
      \item Identification des travaux d'isolation complémentaires recommandés
    \end{itemize}

  \item \textbf{Dimensionnement PAC air/eau} :
    \begin{itemize}
      \item Calcul de la puissance utile nécessaire en régime de base et en point de bivalence
      \item Sélection d'un modèle constructeur (Daikin, Atlantic, Bosch) avec justification
      \item Calcul du COP en condition A7/W35 et du SCOP annuel (méthode bin)
      \item Dimensionnement du ballon tampon (volume, pertes thermiques)
    \end{itemize}

  \item \textbf{Réseau hydraulique} :
    \begin{itemize}
      \item Calcul des débits par circuit (plancher chauffant + radiateurs)
      \item Calcul des pertes de charge (tuyauteries + émetteurs + vannes)
      \item Sélection des circulateurs (courbe pompe/réseau, classe ErP)
      \item Dimensionnement du vase d'expansion (DTU 65.11)
    \end{itemize}

  \item \textbf{Émetteurs} :
    \begin{itemize}
      \item Plancher chauffant RDC : calcul densité de flux, pas de pose, pertes de charge boucles
      \item Radiateurs R+1 : sélection en régime basse température (45/40°C)
    \end{itemize}

  \item \textbf{Aspects économiques} :
    \begin{itemize}
      \item Devis estimatif par poste (€ HT)
      \item Calcul des économies annuelles (fioul → PAC, sur base DJU)
      \item Calcul du TRI avec et sans aides (MaPrimeRénov', éco-PTZ)
    \end{itemize}
\end{enumerate}
\end{methode}

\begin{methode}[Livrables attendus — Note 2 : Soutenance orale]
\textbf{Durée : 20 min exposé + 10 min questions — Semaine 7}

\begin{itemize}
  \item Support visuel : diaporama (10 à 15 slides) + schéma hydraulique A3 imprimé
  \item Chaque membre du groupe présente une partie du dossier
  \item Questions individuelles obligatoires posées par le jury sur les calculs
  \item Critères d'évaluation : clarté technique, rigueur des calculs, justification des choix
\end{itemize}
\end{methode}

% ================================================================
\section{Grille d'évaluation}
% ================================================================

\begin{center}
\begin{tabular}{p{5cm}ccp{5cm}}
  \toprule
  \textbf{Critère} & \textbf{Barème} & \textbf{Note} & \textbf{Indicateurs de réussite} \\
  \midrule
  \multicolumn{4}{l}{\textbf{Note 1 — Dossier technique (20 pts)}} \\
  \midrule
  Bilan thermique & 5 pts & /5 & U corrects, déperditions cohérentes ±10\% \\
  Dim. PAC & 4 pts & /4 & Puissance justifiée, modèle sélectionné \\
  Réseau hydraulique & 4 pts & /4 & Débits, $\Delta P$, circulateur justifiés \\
  Émetteurs & 3 pts & /3 & PCH et radiateurs dimensionnés \\
  Aspect économique & 2 pts & /2 & TRI calculé, aides identifiées \\
  Qualité rédactionnelle & 2 pts & /2 & Plan, syntaxe, figures légendées \\
  \midrule
  \multicolumn{4}{l}{\textbf{Note 2 — Soutenance orale (20 pts)}} \\
  \midrule
  Clarté de l'exposé & 5 pts & /5 & Structure logique, gestion du temps \\
  Maîtrise technique & 8 pts & /8 & Réponses aux questions individuelles \\
  Support visuel & 4 pts & /4 & Schémas clairs, lisibles \\
  Travail d'équipe & 3 pts & /3 & Répartition équilibrée \\
  \bottomrule
\end{tabular}
\end{center}

% ================================================================
\section{Ressources documentaires}
% ================================================================

\begin{itemize}
  \item \textbf{Normes :} NF EN 12831 (calcul besoins chaleur), NF EN 14511 (PAC), DTU 65.11 (sécurités)
  \item \textbf{Réglementation :} RE2020 — arrêté du 4 août 2021, Règlement (UE) 813/2013 (PAC)
  \item \textbf{Aides financières :} MaPrimeRénov' (ANAH), éco-PTZ, TVA 5,5\%
  \item \textbf{Logiciels conseillés :} Pleiades+Comfie (bilan thermique), ClimaWin (PAC)
  \item \textbf{Catalogues constructeurs :} Daikin Altherma 3, Atlantic Extensa+, Buderus Logatherm
  \item \textbf{Cours :} Chapitres 01 (thermique), 02 (hydraulique), 04 (production), 05 (distribution)
\end{itemize}

\finchapitre
